%% kapa: exam=1 round=37 year=2026 period=1 form=A subject=civil_law subject_ko=민법 seq=26 number=26
%% source: Q-Net 감정평가사 기출문제 (cst00309) artlSeq=5258924 file=raw/1008/5258924/2026년 제37회 감정평가사 1차 1교시.pdf
%% conversion: kordoc 4.24.51 -> tools/gen1.py (auto); review=proofread-sonnet; flags=none
\begin{question}{26}
자주점유에 관한 설명으로 옳은 것은? (다툼이 있으면 판례에 따름)
\begin{choices}
  \item 자주점유가 되기 위해서는 점유자가 그 점유의 권원이 있다고 믿어야 한다.
  \item 점유자가 스스로 매매 또는 증여와 같은 자주점유의 권원을 주장하였으나 이것이 인정되지 않으면 이로써 자주점유의 추정은 번복된다.
  \item 매매목적물이 아닌 다른 토지를 인도받아 매매목적물로 믿고 점유한 토지매수인의 점유는 타주점유이다.
  \item 국가나 지방자치단체는 지적공부의 관리주체이므로 부동산을 점유하더라도 자주점유한 것으로 추정되지 않는다.
  \item 타주점유자가 그 명의로 소유권이전등기를 경료한 것만으로는 그 점유를 시킨 자에 대하여 소유의 의사를 표시함으로써 자주점유로 전환되었다고 볼 수 없다.
\end{choices}
\end{question}
\answer{5}
